% sys-first-order-regular-case.tex — \input'd from main.tex
% Status: unapproved draft; agent-written math must be reviewed before integration.
% Sources: HK2017 formula as stated in general-case-algorithm.tex; local
%   behavior classification in research/sys-first-order-local-behavior.md.
% Structure: motivation; support-restricted quadratic problem; generic and
%   feasible candidate definitions; smooth-candidate lemma; sys-gradient
%   theorem; boundary of the generic-feasible argument.

\section{Generic feasible HK gradients}
\label{sec:sys-first-order-generic-feasible}

\begin{unverified}
The Haim--Kislev formula reduces the EHZ capacity of a polytope to finitely
many algebraic candidates. This section treats the easiest case. Suppose that
one candidate is strictly best, that its equality constraints have full rank,
and that its quadratic form is strictly concave along the constraint space.
Then its dwell-time vector varies smoothly when the rows of the polytope move.
The HK quadratic value, the capacity, and the systolic ratio are smooth at such
a point, so their first-order variation is given by an ordinary gradient.

This generic feasible case is not the hard HKO case. It is the baseline
calculation around which the hard cases are organized. The generic-feasible
argument stops at ties, zero dwell-time limits, rank loss, and flat quadratic
directions.

Let
\[
  K(a) = \{x \in \mathbb{R}^4 : \langle a_i,x\rangle \leq 1,\;
    i=1,\ldots,F\}
\]
be a fixed row chart. The rows \(a_i\) are the dual vertices used in the HK
action formula. Fix a nonempty support \(S\subseteq \{1,\ldots,F\}\) and a
cyclic order \(\sigma\) of \(S\), written as
\(\sigma(1),\ldots,\sigma(|S|)\). The support \(S\) records the rows whose
dwell-time coefficients are required to be positive. We use one dwell-time
variable for each position in this cyclic order. Write
\[
  C_{\sigma,S}(a)\beta = e
\]
for the closure and normalization equations
\[
  \sum_{i=1}^{|S|}\beta_i a_{\sigma(i)} = 0,
  \qquad
  \sum_{i=1}^{|S|}\beta_i = 1,
  \qquad
  e=(0,0,0,0,1)^T .
\]
Let \(H_{\sigma,S}(a)\) be the HK action matrix on this ordered support. The
support quadratic is
\[
  Q_{\sigma,S}(\beta,a)
    = \frac12\,\beta^T H_{\sigma,S}(a)\beta .
\]
The support contributes a candidate at \(a\) when the critical point of this
quadratic problem exists and satisfies \(\beta_i>0\) for
\(i=1,\ldots,|S|\).
Coordinates outside \(S\) are omitted; equivalently, their dwell-time
coefficients are set to zero.

The size of \(S\) already determines the generic shape of the constraint
space. Since \(C_{\sigma,S}(a)\) has five rows, supports with
\(|S|\leq 4\) generically have no solution of
\(C_{\sigma,S}(a)\beta=e\). For \(|S|=5\), a rank \(5\) matrix gives one
feasible solution. For \(|S|>5\), a rank \(5\) matrix gives an affine space of
feasible solutions of dimension \(|S|-5\).

We also need a simple volume chamber. Irredundancy of the listed inequalities
means that each listed inequality defines a genuine facet, but it does not by
itself fix the local vertex combinatorics. In this section, the \emph{simple
volume condition at \(a_0\)} means that every vertex of \(K(a_0)\) is cut out
by exactly four listed inequalities, the corresponding \(4\times4\) row
matrices are nonsingular, and every nonincident inequality is strict at every
vertex. These are finitely many strict conditions at \(a_0\), so by continuity
they persist after shrinking to a neighborhood of \(a_0\). On that
neighborhood, the same four-row subsets define the vertices of \(K(a)\), each
vertex is obtained by solving a fixed nonsingular \(4\times4\) linear system
in the row data, and any fixed triangulation of \(K(a_0)\) expresses
\(\operatorname{vol}(K(a))\) as a smooth function of \(a\).

\begin{definition}[HK-generic row parameter]
\label{def:hk-generic-row-parameter}
The row parameter \(a_0\) is \emph{HK-generic in this row chart} if
\(K(a_0)\) satisfies the simple volume condition, and if the following finite
conditions hold for every listed pair \((\sigma,S)\):
\begin{enumerate}
\item if \(|S|\leq4\), then \(C_{\sigma,S}(a_0)\beta=e\) has no solution;
\item if \(|S|\geq5\), then \(\operatorname{rank} C_{\sigma,S}(a_0)=5\);
\item if \(|S|\geq5\), then the restricted quadratic form
  \[
    H_{\sigma,S}(a_0)|_{\ker C_{\sigma,S}(a_0)}
  \]
  has no zero eigenvalue;
\item if \(|S|\geq5\), then the resulting constrained critical point
  \(\beta^*_{\sigma,S}(a_0)\) has no zero coordinate, and
  \(Q_{\sigma,S}(\beta^*_{\sigma,S}(a_0),a_0)\neq0\).
\end{enumerate}
The condition is simultaneous because the list of pairs \((\sigma,S)\) is
finite.
\end{definition}

\begin{definition}[Feasible generic support maximum]
\label{def:feasible-generic-support-maximum}
Let \(a_0\) be HK-generic and let \(|S|\geq 5\). The pair \((\sigma,S)\) is a
\emph{feasible generic support maximum} at \(a_0\) if:
\begin{enumerate}
\item the quadratic form is negative definite on \(\ker C_{\sigma,S}(a_0)\):
  \[
    u^T H_{\sigma,S}(a_0)u < 0
    \qquad
    \text{for all }0\neq u\in \ker C_{\sigma,S}(a_0);
  \]
\item the affine-space maximizer \(\beta^*_{\sigma,S}(a_0)\) is positive on
  the support:
  \[
    \beta^*_{\sigma,S,i}(a_0)>0
    \qquad (i=1,\ldots,|S|);
  \]
\item its objective value is positive:
  \[
    Q_{\sigma,S}(\beta^*_{\sigma,S}(a_0),a_0)>0.
  \]
\end{enumerate}
When \(|S|=5\), the kernel of \(C_{\sigma,S}(a_0)\) is zero under the rank
condition, so the negative-definiteness condition is vacuous and the feasible
\(\beta\) is unique.
\end{definition}

\begin{remark}[Generic versus feasible maximizing]
\label{rem:generic-versus-feasible-maximizing}
The generic conditions are nonvanishing conditions: full rank rather than rank
loss, no zero tangent eigenvalue rather than a flat quadratic direction, no zero
dwell-time coordinate, and no zero value. For a fixed support, these conditions
are open. When the corresponding determinant or numerator is not identically
zero, its nonvanishing locus is dense. Since the HK list is finite, these
genericity tests can be imposed for all supports at once.

Feasible maximizing adds the signs needed for an actual positive support
maximum: \(\beta^*_i>0\) on the support, \(Q>0\), and negative definiteness
rather than mere nondegeneracy of the tangent quadratic form. Thus a generic
support can be discarded because it has a negative dwell-time coordinate, a
negative value, or a positive tangent direction.
\end{remark}

The definition is the standard second-order sufficient condition for an
equality-constrained maximum. To see this, work on a neighborhood where
\(\operatorname{rank} C_{\sigma,S}(a)=5\). Choose a smooth particular solution
\(\beta_0(a)\) of \(C_{\sigma,S}(a)\beta=e\) and a smooth basis matrix
\(N(a)\) for \(\ker C_{\sigma,S}(a)\). Every point in the affine constraint
space has the form
\[
  \beta = \beta_0(a)+N(a)z.
\]
Substitution gives
\[
  Q_{\sigma,S}(\beta_0+Nz,a)
    = \frac12 z^T B(a)z + \ell(a)^Tz + q_0(a),
\]
where
\[
  B(a)=N(a)^T H_{\sigma,S}(a)N(a),
  \qquad
  \ell(a)=N(a)^T H_{\sigma,S}(a)\beta_0(a).
\]
The critical equation on the affine constraint space is
\[
  B(a)z+\ell(a)=0.
\]
If \(B(a)\) is negative definite, then
\[
  z^*(a)=-B(a)^{-1}\ell(a),
  \qquad
  \beta^*_{\sigma,S}(a)=\beta_0(a)+N(a)z^*(a)
\]
exists uniquely and maximizes \(Q_{\sigma,S}\) on the affine constraint space.

\begin{lemma}[Feasible generic support maxima are smooth]
\label{lem:feasible-generic-support-maxima-smooth}
If \((\sigma,S)\) is a feasible generic support maximum at \(a_0\), then the
same support gives a smooth candidate near \(a_0\). More precisely, after
shrinking to a neighborhood of \(a_0\),
\[
  \operatorname{rank} C_{\sigma,S}(a)=5,\qquad
  H_{\sigma,S}(a)|_{\ker C_{\sigma,S}(a)} < 0,
\]
and
\[
  \beta^*_{\sigma,S,i}(a)>0
  \qquad (i=1,\ldots,|S|),
\]
and
\[
  Q^{\mathrm{cand}}_{\sigma,S}(a)
    = Q_{\sigma,S}(\beta^*_{\sigma,S}(a),a)
\]
is a smooth function of \(a\).
\end{lemma}

\begin{proof}
Since \(C_{\sigma,S}(a_0)\) has rank \(5\), some \(5\times 5\) minor is
nonzero at \(a_0\). The same minor remains nonzero for all nearby \(a\), so
the rank stays \(5\). On this patch, \(\beta_0(a)\) and \(N(a)\) may be chosen
smoothly. The matrices \(B(a)\) and vectors \(\ell(a)\) are smooth functions
of \(a\). The eigenvalues of \(B(a)\) vary continuously, so negative
definiteness persists. Thus \(z^*(a)=-B(a)^{-1}\ell(a)\) is smooth. The
strict inequalities \(\beta^*_{\sigma,S,i}(a_0)>0\) persist after shrinking
the neighborhood. Substituting the smooth maximizer into \(Q_{\sigma,S}\)
gives the claimed smooth candidate value.
\end{proof}

\begin{remark}[KKT form]
\label{rem:generic-feasible-kkt-form}
The same proof can be written without choosing \(N(a)\). The equality
constrained critical equations form the saddle-point system
\[
  \begin{pmatrix}
    H_{\sigma,S}(a) & -A_{\sigma,S}(a) & -\mathbf{1} \\
    A_{\sigma,S}(a)^T & 0 & 0 \\
    \mathbf{1}^T & 0 & 0
  \end{pmatrix}
  \begin{pmatrix}
    \beta \\ \lambda \\ \nu
  \end{pmatrix}
  =
  \begin{pmatrix}
    0 \\ 0 \\ 1
\end{pmatrix},
\]
where the columns of \(A_{\sigma,S}(a)^T\) are the selected rows
\(a_{\sigma(i)}\) for \(i=1,\ldots,|S|\). Full constraint rank and negative
definiteness on
\(\ker C_{\sigma,S}(a_0)\) make this system nonsingular. The implicit function
theorem then gives the same smooth candidate value.
\end{remark}

For the next theorem, say that a support candidate is \emph{excluded near}
\(a_0\) if a previously proved pruning result, or an explicit value-gap
calculation on a neighborhood of \(a_0\), shows that it cannot realize the HK
quadratic maximum there. This is a hypothesis of the generic-feasible theorem.
In an application, it must be checked separately.

\begin{theorem}[Smoothness at a unique feasible generic HK candidate]
\label{thm:smooth-sys-feasible-generic-candidate}
Assume that at \(a_0\):
\begin{enumerate}
\item \(K(a_0)\) satisfies the simple volume condition, and
  \(\operatorname{vol}(K(a_0))>0\);
\item one feasible generic support maximum \((\sigma,S)\) has
  \(Q^{\mathrm{cand}}_{\sigma,S}(a_0)\) strictly larger than every other
  feasible generic support maximum that is not excluded near \(a_0\);
\item every support candidate that is not a feasible generic support maximum is
  excluded near \(a_0\).
\end{enumerate}
Then the HK quadratic maximum is locally equal to the smooth candidate value
\(Q^{\mathrm{cand}}_{\sigma,S}(a)\). Consequently
\[
  c_{\mathrm{EHZ}}(K(a))
    = \frac{1}{2 Q^{\mathrm{cand}}_{\sigma,S}(a)}
\]
and
\[
  \operatorname{sys}(K(a))
    = \frac{c_{\mathrm{EHZ}}(K(a))^2}
           {2\,\operatorname{vol}(K(a))}
\]
are smooth near \(a_0\). For every row perturbation \(h\),
\[
  D\operatorname{sys}_{a_0}(h)
    = \nabla \operatorname{sys}(a_0)\cdot h .
\]
\end{theorem}

\begin{proof}
Lemma~\ref{lem:feasible-generic-support-maxima-smooth} gives a smooth candidate value for
\((\sigma,S)\). The strict value gaps at \(a_0\) persist after shrinking the
neighborhood, because only finitely many feasible generic support maxima are
not excluded and their values are continuous on this patch. Thus the same
candidate realizes the HK maximum throughout the neighborhood. The
transformations from \(Q\) to
\(c_{\mathrm{EHZ}}\), and from \(c_{\mathrm{EHZ}}\) to \(\operatorname{sys}\),
are smooth because \(Q^{\mathrm{cand}}_{\sigma,S}(a_0)>0\) and
\(\operatorname{vol}(K(a_0))>0\). The derivative formula is the ordinary
directional derivative of a smooth function.
\end{proof}

\subsection*{Boundary of the generic-feasible argument}

The theorem above is not a theorem about HKO by itself. It is the baseline
generic-feasible calculation. It does not handle:
\begin{itemize}
\item two or more candidates tying for the HK maximum;
\item a positive dwell-time support whose limiting \(\beta\) has a zero
  coordinate;
\item loss of rank in \(C_{\sigma,S}(a)\);
\item a projected Hessian with zero eigenvalues;
\item a non-generic or infeasible candidate that ties the HK maximum.
\end{itemize}
These failures are not cosmetic. They are exactly where the first-order
behavior can become a finite envelope, a one-sided branch problem, or a
singular-support problem. The theorem is still useful because it says
what happens on the open patches around those failures: there the first-order
behavior is computed from ordinary gradients of feasible generic HK candidates.
\end{unverified}
